% GENERATED FILE. Edit canonical lessons, then run npm run generate:books.
% canonical-source-hash: fnv1a64:6b46705ca3afe5e1
% canonical-lessons: TA-C506-uttiyokam, TA-C506-cuyavivaram, TA-C506-amarttu, TA-C506-mulunera, TA-C506-utaviyalar

\chapter{Getting Hired}
\label{ch:ta-getting-hired}
\hlchaptermodality{506}

\begin{chapteropening}
\textbf{By the end of this chapter:} \emph{I can talk about applying for a job and being hired.}

Complete the last lesson of chapter 506: I can talk about applying for a job and being hired.
\end{chapteropening}

\section[\texorpdfstring{\ta{உத்தியோகம்} (uttiyōkam)}{உத்தியோகம் (uttiyōkam)}]{\ta{உத்தியோகம்} (uttiyōkam) — a salaried job; a post}

\label{lesson:TA-C506-uttiyokam}



\begin{quote}
Before the new one: say the Tamil for a complaint, then the Tamil for an athlete.
\end{quote}

\subsection*{You'll want to know: \ta{உத்தியோகம்}}
\textbf{\ta{உத்தியோகம்}} — \emph{uttiyōkam} — \textquotedblleft{}a salaried job; a post\textquotedblright{}.

\begin{quote}
\ta{அவருக்கு} \ta{வங்கியில்} \ta{உத்தியோகம்} \ta{கிடைத்தது}.
\end{quote}

\begin{quote}
\emph{avarukku vaṅkiyil uttiyōkam kiṭaittatu.}
\end{quote}

\begin{quote}
He got a job in a bank.
\end{quote}

A more formal word than \ta{வேலை} (\emph{vēlai}): \ta{உத்தியோகம்} suggests a salaried post, often in an office.

\begin{cousinweb}
\textbf{\ta{உத்தியோகம்}} (\emph{uttiyōkam}) is borrowed from Sanskrit \emph{udyoga}, \textquotedblleft{}active preparation, exertion\textquotedblright{} (CDIAL 2055). Tamil wrote the Sanskrit cluster \emph{dy} as \ta{த்திய} and gave the noun the ending -\ta{அம்}.
\end{cousinweb}

\begin{grammarlens}[title={What you've built}]
One more word for describing what you lived and felt.
\end{grammarlens}

\subsection*{Your turn}
\begin{itemize}
\raggedright
  \item \emph{Say it:} \emph{uttiyōkam}
  \item \emph{Say it:} \emph{uttiyōkam}, once more
  \item \emph{Say it:} \emph{avarukku vaṅkiyil uttiyōkam kiṭaittatu.} — He got a job in a bank.
  \item \emph{Recall:} say the Tamil for a complaint, then the Tamil for an athlete, then say \emph{uttiyōkam} again
\end{itemize}

\begin{tcolorbox}[breakable,colback=teal!4,colframe=teal!35!black,title={Before you move on}]
What does \ta{உத்தியோகம்} mean? (\textquotedblleft{}A salaried job; a post\textquotedblright{}.) Say it once more.
\end{tcolorbox}
\section[\texorpdfstring{\ta{சுயவிவரம்} (cuyavivaram)}{சுயவிவரம் (cuyavivaram)}]{\ta{சுயவிவரம்} (cuyavivaram) — a CV; a profile}

\label{lesson:TA-C506-cuyavivaram}



\begin{quote}
Before the new one: say the Tamil for an athlete, then the Tamil for a salaried job.
\end{quote}

\subsection*{You'll want to know: \ta{சுயவிவரம்}}
\textbf{\ta{சுயவிவரம்}} — \emph{cuyavivaram} — \textquotedblleft{}a CV; a profile\textquotedblright{}.

\begin{quote}
\ta{நான்} \ta{என்} \ta{சுயவிவரத்தை} \ta{மின்னஞ்சலில்} \ta{அனுப்பினேன்}.
\end{quote}

\begin{quote}
\emph{nāṉ eṉ cuyavivarattai miṉṉañcalil aṉuppiṉēṉ.}
\end{quote}

\begin{quote}
I sent my CV by email.
\end{quote}

A noun in -\ta{அம்} takes -\ta{அத்து}- before an ending: \ta{சுயவிவரத்தை}, \ta{சுயவிவரத்தில்}.

\begin{cousinweb}
\textbf{\ta{சுயவிவரம்}} (\emph{cuyavivaram}) is a modern compound of \ta{சுய} (\emph{cuya}), \textquotedblleft{}one's own\textquotedblright{}, from Sanskrit \emph{sva}, and \ta{விவரம்} (\emph{vivaram}), \textquotedblleft{}details\textquotedblright{}, borrowed from Sanskrit \emph{vivaraṇa}, \textquotedblleft{}explanation\textquotedblright{} (Monier-Williams).
\end{cousinweb}

\begin{grammarlens}[title={What you've built}]
One more word for describing what you lived and felt.
\end{grammarlens}

\subsection*{Your turn}
\begin{itemize}
\raggedright
  \item \emph{Say it:} \emph{cuyavivaram}
  \item \emph{Say it:} \emph{cuyavivaram}, once more
  \item \emph{Say it:} \emph{nāṉ eṉ cuyavivarattai miṉṉañcalil aṉuppiṉēṉ.} — I sent my CV by email.
  \item \emph{Recall:} say the Tamil for an athlete, then the Tamil for a salaried job, then say \emph{cuyavivaram} again
\end{itemize}

\begin{tcolorbox}[breakable,colback=teal!4,colframe=teal!35!black,title={Before you move on}]
What does \ta{சுயவிவரம்} mean? (\textquotedblleft{}A CV; a profile\textquotedblright{}.) Say it once more.
\end{tcolorbox}
\section[\texorpdfstring{\ta{அமர்த்து} (amarttu)}{அமர்த்து (amarttu)}]{\ta{அமர்த்து} (amarttu) — to hire; to seat}

\label{lesson:TA-C506-amarttu}



\begin{quote}
Before the new one: say the Tamil for a salaried job, then the Tamil for a CV.
\end{quote}

\subsection*{You'll want to know: \ta{அமர்த்து}}
\textbf{\ta{அமர்த்து}} — \emph{amarttu} — \textquotedblleft{}to hire; to seat\textquotedblright{}.

\begin{quote}
\ta{அந்த} \ta{நிறுவனம்} \ta{புதிய} \ta{ஓட்டுநரை} \ta{அமர்த்தியது}.
\end{quote}

\begin{quote}
\emph{anta niṟuvaṉam putiya ōṭṭunarai amarttiyatu.}
\end{quote}

\begin{quote}
That company hired a new driver.
\end{quote}

Past (I …): \ta{அமர்த்தினேன்}, \emph{amarttiṉēṉ}.

The person hired takes -\ta{ஐ}: \ta{ஓட்டுநரை} \ta{அமர்த்து}. The full phrase is \ta{வேலைக்கு} \ta{அமர்த்து} (\emph{vēlaikku amarttu}), to take someone on for work.

\begin{cousinweb}
\textbf{\ta{அமர்த்து}} (\emph{amarttu}) is native Dravidian: DEDR 161 lists \ta{அமர்} (\emph{amar}), \textquotedblleft{}to settle, to rest\textquotedblright{}, and beside it \ta{அமர்த்து}, \textquotedblleft{}to engage a servant, to establish\textquotedblright{}. The strong verb makes someone settle into a place or a post.
\end{cousinweb}

\begin{grammarlens}[title={What you've built}]
One more word for describing what you lived and felt.
\end{grammarlens}

\subsection*{Your turn}
\begin{itemize}
\raggedright
  \item \emph{Say it:} \emph{amarttu}
  \item \emph{Say it:} \emph{amarttu}, once more
  \item \emph{Say it:} \emph{anta niṟuvaṉam putiya ōṭṭunarai amarttiyatu.} — That company hired a new driver.
  \item \emph{Recall:} say the Tamil for a salaried job, then the Tamil for a CV, then say \emph{amarttu} again
\end{itemize}

\begin{tcolorbox}[breakable,colback=teal!4,colframe=teal!35!black,title={Before you move on}]
What does \ta{அமர்த்து} mean? (\textquotedblleft{}To hire; to seat\textquotedblright{}.) Say it once more.
\end{tcolorbox}
\section[\texorpdfstring{\ta{முழுநேர} (muḻunēra)}{முழுநேர (muḻunēra)}]{\ta{முழுநேர} (muḻunēra) — full-time}

\label{lesson:TA-C506-mulunera}



\begin{quote}
Before the new one: say the Tamil for a CV, then the Tamil for to hire.
\end{quote}

\subsection*{You'll want to know: \ta{முழுநேர}}
\textbf{\ta{முழுநேர}} — \emph{muḻunēra} — \textquotedblleft{}full-time\textquotedblright{}.

\begin{quote}
\ta{என்} \ta{அக்கா} \ta{ஒரு} \ta{முழுநேர} \ta{வேலை} \ta{தேடுகிறாள்}.
\end{quote}

\begin{quote}
\emph{eṉ akkā oru muḻunēra vēlai tēṭukiṟāḷ.}
\end{quote}

\begin{quote}
My elder sister is looking for a full-time job.
\end{quote}

It stands before a noun, like \ta{பெரிய} (\emph{periya}): \ta{முழுநேர} \ta{வேலை}, \ta{முழுநேர} \ta{மாணவன்}.

\begin{cousinweb}
\textbf{\ta{முழுநேர}} (\emph{muḻunēra}) is a native compound: \ta{முழு} (\emph{muḻu}), \textquotedblleft{}whole, entire\textquotedblright{} (DEDR 4992), and \ta{நேரம்} (\emph{nēram}), \textquotedblleft{}time\textquotedblright{} (DEDR 3774), in the shape a noun takes before another noun.
\end{cousinweb}

\begin{grammarlens}[title={What you've built}]
One more word for describing what you lived and felt.
\end{grammarlens}

\subsection*{Your turn}
\begin{itemize}
\raggedright
  \item \emph{Say it:} \emph{muḻunēra}
  \item \emph{Say it:} \emph{muḻunēra}, once more
  \item \emph{Say it:} \emph{eṉ akkā oru muḻunēra vēlai tēṭukiṟāḷ.} — My elder sister is looking for a full-time job.
  \item \emph{Recall:} say the Tamil for a CV, then the Tamil for to hire, then say \emph{muḻunēra} again
\end{itemize}

\begin{tcolorbox}[breakable,colback=teal!4,colframe=teal!35!black,title={Before you move on}]
What does \ta{முழுநேர} mean? (\textquotedblleft{}Full-time\textquotedblright{}.) Say it once more.
\end{tcolorbox}
\section[\texorpdfstring{\ta{உதவியாளர்} (utaviyāḷar)}{உதவியாளர் (utaviyāḷar)}]{\ta{உதவியாளர்} (utaviyāḷar) — an assistant}

\label{lesson:TA-C506-utaviyalar}



\begin{quote}
Before the new one: say the Tamil for to hire, then the Tamil for full-time.
\end{quote}

\subsection*{You'll want to know: \ta{உதவியாளர்}}
\textbf{\ta{உதவியாளர்}} — \emph{utaviyāḷar} — \textquotedblleft{}an assistant\textquotedblright{}.

\begin{quote}
\ta{மருத்துவரின்} \ta{உதவியாளர்} \ta{என்} \ta{பெயரை} \ta{அழைத்தார்}.
\end{quote}

\begin{quote}
\emph{maruttuvariṉ utaviyāḷar eṉ peyarai aḻaittār.}
\end{quote}

\begin{quote}
The doctor's assistant called out my name.
\end{quote}

Plural: \ta{உதவியாளர்கள்}, \emph{utaviyāḷarkaḷ}.

The ending -\ta{ஆளர்} makes a noun for \textquotedblleft{}one who …\textquotedblright{}: \ta{உதவி} \textquotedblleft{}help\textquotedblright{} + \ta{ஆளர்}. Compare \ta{காசாளர்} (\emph{kācāḷar}) and \ta{பொறியாளர்} (\emph{poṟiyāḷar}).

\begin{cousinweb}
\textbf{\ta{உதவியாளர்}} (\emph{utaviyāḷar}) is built from native parts: \ta{உதவி} (\emph{utavi}), \textquotedblleft{}help\textquotedblright{}, from \ta{உதவு}, \textquotedblleft{}to help\textquotedblright{} (DEDR 609), and the ending -\ta{ஆளர்}, \textquotedblleft{}one who does it\textquotedblright{}.
\end{cousinweb}

\begin{grammarlens}[title={What you've built}]
One more word for describing what you lived and felt.
\end{grammarlens}

\subsection*{Your turn}
\begin{itemize}
\raggedright
  \item \emph{Say it:} \emph{utaviyāḷar}
  \item \emph{Say it:} \emph{utaviyāḷar}, once more
  \item \emph{Say it:} \emph{maruttuvariṉ utaviyāḷar eṉ peyarai aḻaittār.} — The doctor's assistant called out my name.
  \item \emph{Recall:} say the Tamil for to hire, then the Tamil for full-time, then say \emph{utaviyāḷar} again
\end{itemize}

\begin{tcolorbox}[breakable,colback=teal!4,colframe=teal!35!black,title={Before you move on}]
What does \ta{உதவியாளர்} mean? (\textquotedblleft{}An assistant\textquotedblright{}.) Say it once more.
\end{tcolorbox}
